\chapter{Synthèse des résultats}\label{ch:synthese}

\section{Un corpus défini et contrôlé}

Le catalogue conserve les \st{nA} notices de la version précédente, avec leur numérotation. La confrontation des témoins a conduit à reconnaître trois paires de notices qui transmettent le même monument, un fragment à joindre à une autre notice et un texte extérieur à la période. Les analyses portent sur \st{nB} textes (ensemble B), dont \st{nC} textes funéraires (ensemble C) ; les deux dépositions successives d’une même plaque (\nnums{313} et 314) forment un seul monument, ce qui donne \st{nD} unités matérielles (chapitre~\ref{ch:corpus}).

Le texte de chaque notice est désormais confronté à un second témoin indépendant de l’EDCS. L’EDB a été appariée à \st{edbApp} notices ; pour \st{edbConc} d’entre elles, auxquelles s’ajoutent \st{edbConcRel} cas après relecture de la fiche détaillée, les deux lectures concordent ; pour \st{edbExam}, elles présentent des écarts mineurs ; pour \st{edbDiv}, une divergence forte ou un appariement douteux impose le retour à l’édition imprimée (chapitre~\ref{ch:transmission}). La seconde base apporte en outre ce que la première ignorait : le support, ses dimensions, la technique d’exécution, les signes figurés, le rattachement topographique, le lieu de conservation et les reproductions disponibles.

\section{Des années rapportées à leur fondement}

Chaque année de classement est rapportée à ce qui la fonde (chapitre~\ref{ch:fondement}). Dans l’ensemble C, l’année de \st{datLu} textes repose sur un nom consulaire lu, celle de \st{datPartiel} sur une lecture partielle, celle de \st{datInit} sur des initiales et celle de \st{datRest} sur un consulat restitué. Huit formules de date ne sont transmises que par l’EDB ; quatre années reposent sur une restitution conjecturale, quatre sur la seule base. Dix datations sont contestées, cinq font l’objet d’une réserve. Aucune année n’a été modifiée par vraisemblance ; les sept corrections apportées à l’année de la base reposent sur le texte.

Les contrôles calendaires éprouvent la cohérence des dates du jour : sur dix-huit mentions du jour de la semaine, dix concordent avec le calcul julien (dont une si l’on retient l’année 398 pour le \nnum{428}), cinq présentent un écart d’un jour qui reste inexpliqué, une est discordante et deux ne sont pas vérifiables ; et les deux mentions de l’âge de la lune sont compatibles avec les années proposées (chapitre~\ref{ch:calendrier}).

\section{Le formulaire, les signes et la langue}

La formule \emph{in pace} figure dans \st{ip_tot} textes funéraires (\st{ip_pct}~\%), dont près d’un quart seulement par restitution. Le noyau du formulaire (\emph{vixit}, \emph{depositus}, \emph{in pace}) est stable sur toute la période ; son recul apparent dans les textes tardifs tient en partie à leur état de conservation. Deux évolutions résistent mieux à ce biais : le recul des termes conjugaux et l’apparition, dans les dernières décennies, de \emph{plus minus}, des mentions d’emplacement et des actes d’achat (chapitre~\ref{ch:formulaire}).

Le chrisme et la croix sont décrits sur \st{nSignesC} textes funéraires, soit \st{pSignesC}~\% des textes appariés. Leur part paraît décroître au cours de la période ; ils sont plus fréquents sur les textes qui portent \emph{in pace} (chapitre~\ref{ch:supports}). La langue des épitaphes présente les traits graphiques du latin tardif ; l’accusatif de durée et l’ablatif coexistent, et la plupart des formes développées relèvent de l’éditeur (chapitre~\ref{ch:langue}).

\section{Les personnes}

L’ensemble C nomme \st{nPersAge} personnes dont l’âge est indiqué, après dix corrections de la table des personnes. L’âge médian est de \st{ageMed} ans ; l’attraction des âges ronds est forte (indice de Whipple : \st{whipple}), et l’expression \emph{plus minus} reconnaît l’approximation. Les épitaphes documentent le couple, la commémoration des enfants, le clergé marié et les lecteurs des titres, quelques clarissimes et quelques étrangers ; elles ne donnent pas une image de la population chrétienne de Rome (chapitres~\ref{ch:personnes} à \ref{ch:relations}).

\section{Le vocabulaire religieux}

Le langage doctrinal est rare. Selon la règle adoptée, \st{nProfilLarge}, \st{nProfilRestr} ou \st{nProfilStrict} textes désignent Dieu, le Christ, l’Esprit ou la Trinité par un mot lu ou abrégé (2,8, 2,4 ou 1,7~\%) ; \st{nProfilAvecRest} avec les restitutions. La Trinité est mentionnée une fois (\nnum{474}), l’Esprit saint au plus une fois et sur restitution (\nnum{249}), la résurrection jamais de manière assurée. Une seule demande d’intercession est assurée (\nnum{214}). L’identité chrétienne s’exprime surtout par le repos et la paix, par les signes, et par la qualification des personnes : sept néophytes, un catéchumène, des \emph{fideles}, des \emph{ancillae dei}, des vierges (chapitres~\ref{ch:typologie} à \ref{ch:textes_decisifs}).

\section{Les sépultures}

Quatorze textes enregistrent l’acquisition d’un emplacement, tous postérieurs à 369. Les vendeurs nommés sont des \emph{fossores} ; l’intervention du clergé n’est attestée que par des restitutions. La sépulture près des saints est présentée comme un privilège rare (\nnum{232}) (chapitre~\ref{ch:sepultures}).
